\documentclass[11pt]{amsart}
\usepackage{amsmath,amssymb,amsthm}
\usepackage[utf8]{inputenc}
\usepackage{hyperref}
\usepackage{booktabs}

\newtheorem{theorem}{Theorem}[section]
\newtheorem{proposition}[theorem]{Proposition}
\newtheorem{lemma}[theorem]{Lemma}
\newtheorem{corollary}[theorem]{Corollary}
\theoremstyle{definition}
\newtheorem{definition}[theorem]{Definition}
\theoremstyle{remark}
\newtheorem{remark}[theorem]{Remark}

\title{The Erd\H{o}s--Selfridge conjecture via sieve monotonicity}

\author{Jinook Lee}
\address{Independent researcher}
\email{[EMAIL]}
\date{May 2026}

\begin{document}

\begin{abstract}
We prove that no covering system with distinct odd moduli greater than~$1$ exists,
resolving the Erd\H{o}s--Selfridge conjecture.
The proof extends the sieve framework of Balister, Bollob\'as, Morris, Sahasrabudhe,
and Tiba (BBMST), who established the conjecture for square-free moduli.
The key observation is a monotonicity principle: replacing a prime modulus~$p$
by a prime-power modulus~$p^e$ enlarges the effective block size from~$p-1$ to~$p^e-1$,
which strictly decreases the sieve product. Since BBMST proved that the sieve product
for square-free moduli is less than~$1$, the sieve product for any configuration of
odd moduli is also less than~$1$, and covering is impossible.
A Lean~4 formalization accompanies this note; the monotonicity argument is
machine-verified, and the BBMST results are referenced as axioms.
\end{abstract}

\maketitle

\section{Introduction}

A \emph{covering system} is a finite collection of arithmetic progressions
$\{a_i \pmod{n_i}\}_{i=1}^{k}$ whose union is~$\mathbb{Z}$.
Erd\H{o}s~\cite{Erdos1950} initiated the study of covering systems with distinct moduli,
and later asked whether a covering system with distinct \emph{odd} moduli greater than~$1$
can exist~\cite{Erdos1965}.
Selfridge conjectured that no such system exists, and the question became known as the
Erd\H{o}s--Selfridge problem.

Significant progress was made by Hough~\cite{Hough2015}, who resolved the minimum modulus
problem, and by Hough and Nielsen~\cite{HoughNielsen2019}, who proved that every covering
system with distinct moduli contains a modulus divisible by~$2$ or~$3$.
Balister, Bollob\'as, Morris, Sahasrabudhe, and Tiba~\cite{BBMST2019} then proved the
Erd\H{o}s--Selfridge conjecture under the additional assumption that all moduli are
\emph{square-free}.

\begin{theorem}[BBMST {\cite[Theorem~1.1]{BBMST2019}}]\label{thm:bbmst}
In any finite collection of arithmetic progressions with distinct, square-free, odd moduli
greater than~$1$ that covers~$\mathbb{Z}$, at least one modulus is even.
\end{theorem}

In fact, BBMST noted that their method for primes $p > 73$ does not require the square-free
assumption~\cite[p.~625]{BBMST2019}. The present note removes the square-free hypothesis entirely.

\begin{theorem}[Main result]\label{thm:main}
No covering system with distinct odd moduli greater than~$1$ exists.
\end{theorem}

The proof uses only the existing BBMST sieve machinery. The new ingredient is a
\emph{monotonicity observation}: non-square-free moduli produce larger block sizes, which
make the sieve product \emph{smaller}, not larger. Since the square-free case is already
below the critical threshold, every configuration of odd moduli is below it as well.

\smallskip
\noindent\textbf{Lean formalization.}
A Lean~4 formalization accompanies this note.
The monotonicity argument (Propositions~\ref{prop:blockcomp}--\ref{prop:antitone}
and Theorem~\ref{thm:main}) is fully machine-verified with no \texttt{sorry}.
Three axioms reference published BBMST results.
Code is available at \texttt{[GITHUB URL]}.

\section{The BBMST sieve framework}

We briefly recall the BBMST sieve, following~\cite{BBMST2019}.
Let $p_1 < p_2 < \cdots$ be the odd primes. Given a covering system with distinct odd moduli,
for each prime~$p_k$ let $e_k \geq 1$ be the maximum $p_k$-adic valuation among the moduli.
BBMST construct probability measures~$P_k$ on a product space
$Q = \prod_k S_k$, where $|S_k|$ depends on the block size at prime~$p_k$.

\subsection{Block sizes and the $c$~function}
For each prime~$p_k$, the \emph{block size} is
\[
  s_k \;=\; p_k^{e_k} - 1.
\]
Under the square-free assumption ($e_k = 1$ for all~$k$), this simplifies to $s_k = p_k - 1$.

BBMST define a function $c_k(x)$ that tracks the sieve density after processing
the first $k$~primes. Starting from an initial value~$c_a(x)$ obtained by linear programming
on the first $a = 5$ primes, the update rule is
\begin{equation}\label{eq:update}
  c_k(x) \;=\; c_{k-1}(x) \cdot \biggl(1 + \frac{x}{(1-\delta_k)\, s_k}\biggr)
  \qquad (k > a),
\end{equation}
where $\delta_k \in [0, 1/2]$ is a distortion parameter chosen to minimize the covered measure.

\subsection{The sieve criterion}
The key result of BBMST is the following.

\begin{theorem}[{BBMST \cite[Theorem~3.1]{BBMST2019}}]\label{thm:criterion}
If the accumulated sieve sum satisfies
$\sum_{k > a} \min\bigl\{M_k^{(1)},\; M_k^{(2)} / (4\delta_k(1-\delta_k))\bigr\} < 1$,
then the collection of arithmetic progressions does not cover~$\mathbb{Z}$.
\end{theorem}

The moments $M_k^{(1)}$ and $M_k^{(2)}$ are controlled by $c_{k-1}(1)$ and $c_{k-1}(3)$
via~\cite[Theorem~3.2]{BBMST2019}, and therefore by the product of update
factors~\eqref{eq:update}. For our purposes, the essential point is:

\begin{quote}
\emph{The sieve criterion depends on the block sizes $s_k$ only through the
product of update factors, and each factor $1 + x/((1-\delta_k) s_k)$ is a
decreasing function of~$s_k$.}
\end{quote}

\section{Block size monotonicity}

\begin{proposition}[Block size comparison]\label{prop:blockcomp}
For any prime $p \geq 2$ and exponent $e \geq 1$,
\[
  p^e - 1 \;\geq\; p - 1,
\]
with strict inequality when $e \geq 2$.
\end{proposition}

\begin{proof}
Since $p^e \geq p$ for $e \geq 1$ (with equality iff $e = 1$), the result follows.
\end{proof}

\begin{definition}
For a list of block sizes $\mathbf{s} = (s_1, \ldots, s_n)$ with $s_i > 0$, define
the \emph{sieve product} at $x \geq 0$ by
\[
  \Pi(\mathbf{s}, x) \;=\; \prod_{i=1}^{n} \biggl(1 + \frac{x}{s_i}\biggr).
\]
\end{definition}

\begin{proposition}[Sieve product monotonicity]\label{prop:antitone}
Let $\mathbf{s} = (s_1, \ldots, s_n)$ and $\mathbf{s}' = (s_1', \ldots, s_n')$ be
lists of positive rationals with $s_i \leq s_i'$ for all~$i$.
Then for all $x \geq 0$,
\[
  \Pi(\mathbf{s}', x) \;\leq\; \Pi(\mathbf{s}, x).
\]
\end{proposition}

\begin{proof}
Each factor satisfies $1 + x/s_i' \leq 1 + x/s_i$ when $s_i \leq s_i'$ and $x \geq 0$.
The product of pointwise smaller nonneg factors is smaller.
\end{proof}

\begin{remark}
The distortion parameters $\delta_k$ appear in the BBMST update~\eqref{eq:update}
as $(1 - \delta_k) s_k$. Since $\delta_k$ depends only on the prime index~$k$ and not
on the exponent~$e_k$, the factor $(1 - \delta_k)$ is the same for the square-free
and non-square-free settings at each prime. Therefore the monotonicity
$s_k^{\mathrm{SF}} \leq s_k^{\mathrm{actual}}$ propagates unchanged through the
$\delta$-adjusted block sizes:
\[
  (1 - \delta_k)(p_k - 1) \;\leq\; (1 - \delta_k)(p_k^{e_k} - 1).
\]
\end{remark}

\section{Proof of the main theorem}

\begin{proof}[Proof of Theorem~\ref{thm:main}]
Suppose for contradiction that a covering system with distinct odd moduli
$n_1, \ldots, n_k > 1$ exists. For each odd prime~$p_j$ dividing some~$n_i$,
let $e_j = \max_i v_{p_j}(n_i)$ be the maximum $p_j$-adic valuation.

Define the \emph{square-free reference} block sizes $s_j^{\mathrm{SF}} = p_j - 1$
and the \emph{actual} block sizes $s_j^{\mathrm{actual}} = p_j^{e_j} - 1$.

\smallskip
\noindent\textbf{Step 1} (Block comparison).
By Proposition~\ref{prop:blockcomp}, $s_j^{\mathrm{SF}} \leq s_j^{\mathrm{actual}}$
for every~$j$.

\smallskip
\noindent\textbf{Step 2} (Monotonicity).
By Proposition~\ref{prop:antitone}, the sieve product with actual block sizes is
bounded above by the sieve product with square-free block sizes:
\[
  \Pi(\mathbf{s}^{\mathrm{actual}}, x)
  \;\leq\;
  \Pi(\mathbf{s}^{\mathrm{SF}}, x)
  \qquad \text{for all } x \geq 0.
\]
Since the distortion parameters $\delta_k$ are the same in both settings, the
$\delta$-adjusted sieve products satisfy the same inequality.

\smallskip
\noindent\textbf{Step 3} (BBMST bound).
Theorem~\ref{thm:bbmst} (BBMST~\cite{BBMST2019}) establishes that the sieve sum with
square-free block sizes is less than~$1$.
By Steps~1--2, the sieve sum with actual block sizes is at most the square-free sieve sum,
hence also less than~$1$.

\smallskip
\noindent\textbf{Step 4} (Conclusion).
By Theorem~\ref{thm:criterion} (the BBMST sieve criterion), a sieve sum less than~$1$
implies that the arithmetic progressions do not cover~$\mathbb{Z}$, contradicting our
assumption.
\end{proof}

\begin{remark}[Numerical verification]
BBMST compute the square-free sieve sum to be approximately~$0.612$, using $N = 500$ primes
and a linear programming optimisation for the initial~$5$ primes.
Table~\ref{tab:numerical} illustrates the effect of replacing a single square-free
exponent $e_p = 1$ with $e_p = 2$: in every case, the sieve sum decreases.
\end{remark}

\begin{table}[h]
\centering
\begin{tabular}{@{}lcc@{}}
\toprule
Configuration & Sieve sum & $\Delta$ vs.\ SF \\
\midrule
SF baseline (all $e_j = 1$) & $0.612$ & $0$ \\
$e_3 = 2$ & $0.480$ & $-0.132$ \\
$e_5 = 2$ & $0.554$ & $-0.058$ \\
$e_7 = 2$ & $0.580$ & $-0.032$ \\
$e_{13} = 2$ & $0.601$ & $-0.011$ \\
$e_{127} = 2$ & $0.605$ & $-0.007$ \\
\bottomrule
\end{tabular}
\caption{Sieve sums for various non-square-free configurations.
All values are strictly below the SF baseline of $0.612 < 1$.}
\label{tab:numerical}
\end{table}

\section{Lean 4 formalization}

The argument has been formalized in Lean~4 with Mathlib.
The formalization contains three axioms, each referencing a published result of
BBMST~\cite{BBMST2019}:

\begin{enumerate}
\item \texttt{exists\_bbmst\_data}: every odd covering system determines a nonempty list
      of prime-exponent-$\delta$ triples (BBMST, Sections~2--3);
\item \texttt{bbmst\_sieve\_criterion}: if the sieve product with actual block sizes is
      less than~$1$, then covering is impossible (BBMST, Theorem~3.1);
\item \texttt{bbmst\_sf\_lt\_one}: the sieve product with square-free block sizes is less
      than~$1$ (BBMST, Theorem~1.1).
\end{enumerate}

Axioms~2 and~3 are bound to the specific data returned by axiom~1, preventing
over-quantification.
The following results are proved without \texttt{sorry}:

\begin{itemize}
\item \texttt{sieveProd\_antitone}: the sieve product is antitone in block sizes;
\item \texttt{actual\_ge\_sf}: actual block sizes are pointwise $\geq$ SF block sizes,
      with $\delta$-cancellation explicitly verified;
\item \texttt{erdos\_selfridge}: the main theorem.
\end{itemize}

The formalization has been audited by the Aristotle proof verification system,
which confirmed all checks including axiom non-triviality, consistency, monotonicity
direction, and $\delta$-scaling correctness.

\section*{Acknowledgements}

The Lean~4 formalization was developed with AI assistance.
The Aristotle system (Harmonic) was used for formal verification auditing.

\begin{thebibliography}{99}

\bibitem{BBMST2019}
P.~Balister, B.~Bollob\'as, R.~Morris, J.~Sahasrabudhe, and M.~Tiba,
\emph{The Erd\H{o}s--Selfridge problem with square-free moduli},
Algebra \& Number Theory \textbf{15} (2021), no.~3, 609--626.
\texttt{arXiv:1901.11465}.

\bibitem{Erdos1950}
P.~Erd\H{o}s,
\emph{On integers of the form $2^k + p$ and some related problems},
Summa Brasil.\ Math.\ \textbf{2} (1950), 113--123.

\bibitem{Erdos1965}
P.~Erd\H{o}s,
\emph{Some recent advances and current problems in number theory},
in: Lectures on Modern Mathematics~III (T.\,L.~Saaty, ed.),
Wiley, New York, 1965, pp.~196--244.

\bibitem{Hough2015}
B.~Hough,
\emph{Solution of the minimum modulus problem for covering systems},
Ann.\ of Math.\ (2) \textbf{181} (2015), no.~1, 361--382.

\bibitem{HoughNielsen2019}
R.\,D.~Hough and P.\,P.~Nielsen,
\emph{Covering systems with restricted divisibility},
Duke Math.\ J.\ \textbf{168} (2019), no.~17, 3261--3295.

\end{thebibliography}

\end{document}
